% Options for packages loaded elsewhere
\PassOptionsToPackage{unicode}{hyperref}
\PassOptionsToPackage{hyphens}{url}
\documentclass[
]{article}
\usepackage{xcolor}
\usepackage{amsmath,amssymb}
\setcounter{secnumdepth}{-\maxdimen} % remove section numbering
\usepackage{iftex}
\ifPDFTeX
  \usepackage[T1]{fontenc}
  \usepackage[utf8]{inputenc}
  \usepackage{textcomp} % provide euro and other symbols
\else % if luatex or xetex
  \usepackage{unicode-math} % this also loads fontspec
  \defaultfontfeatures{Scale=MatchLowercase}
  \defaultfontfeatures[\rmfamily]{Ligatures=TeX,Scale=1}
\fi
\usepackage{lmodern}
\ifPDFTeX\else
  % xetex/luatex font selection
\fi
% Use upquote if available, for straight quotes in verbatim environments
\IfFileExists{upquote.sty}{\usepackage{upquote}}{}
\IfFileExists{microtype.sty}{% use microtype if available
  \usepackage[]{microtype}
  \UseMicrotypeSet[protrusion]{basicmath} % disable protrusion for tt fonts
}{}
\makeatletter
\@ifundefined{KOMAClassName}{% if non-KOMA class
  \IfFileExists{parskip.sty}{%
    \usepackage{parskip}
  }{% else
    \setlength{\parindent}{0pt}
    \setlength{\parskip}{6pt plus 2pt minus 1pt}}
}{% if KOMA class
  \KOMAoptions{parskip=half}}
\makeatother
\usepackage{longtable,booktabs,array}
\newcounter{none} % for unnumbered tables
\usepackage{calc} % for calculating minipage widths
% Correct order of tables after \paragraph or \subparagraph
\usepackage{etoolbox}
\makeatletter
\patchcmd\longtable{\par}{\if@noskipsec\mbox{}\fi\par}{}{}
\makeatother
% Allow footnotes in longtable head/foot
\IfFileExists{footnotehyper.sty}{\usepackage{footnotehyper}}{\usepackage{footnote}}
\makesavenoteenv{longtable}
\setlength{\emergencystretch}{3em} % prevent overfull lines
\providecommand{\tightlist}{%
  \setlength{\itemsep}{0pt}\setlength{\parskip}{0pt}}
\usepackage{bookmark}
\IfFileExists{xurl.sty}{\usepackage{xurl}}{} % add URL line breaks if available
\urlstyle{same}
\hypersetup{
  hidelinks,
  pdfcreator={LaTeX via pandoc}}

\author{}
\date{}

\begin{document}

\section{Read the Room, Not the Files: MCP Filesystem Intelligence for
Privacy-Respecting Compliance
Scanning}\label{read-the-room-not-the-files-mcp-filesystem-intelligence-for-privacy-respecting-compliance-scanning}

\textbf{Adhithya Rajasekaran} and \textbf{Lakshman Shanmugam}

rajasekaran.adhit@gmail.com \textbar{} slakshman2004@gmail.com

\begin{center}\rule{0.5\linewidth}{0.5pt}\end{center}

\subsection{Abstract}\label{abstract}

Enterprise compliance scanners inspect every file on every endpoint to
find sensitive data. On a developer workstation with 30,000 files, fewer
than 500 typically contain anything compliance-relevant. Current Data
Loss Prevention (DLP) systems have no contextual understanding of
filesystem structure: a real API key in a production configuration is
treated identically to a placeholder in a test file, producing
false-positive rates exceeding 80\% in practice. This paper presents a
three-phase scanning architecture built on the Model Context Protocol
(MCP). In Phase 1, an MCP server exposes filesystem metadata ---
directory structure, filenames, extensions, permissions, and file sizes
--- as structured tool calls. No file contents are read. In Phase 2, a
risk engine combines keyword signals, graph centrality, and semantic
directory inference to score every file and select the top 1--5\% for
inspection. In Phase 3, only these targeted files undergo content-based
pattern matching with compliance validation. The system was implemented
as a FastAPI application converted to an MCP server via
\texttt{fastapi\_mcp}, exposing 15 tools for filesystem inspection, risk
scoring, targeted scanning, and analyst feedback. Evaluation across six
real repositories (235 to 29,339 files) demonstrates consistent 95\%
scan reduction with mean precision of 1.00 and mean recall of 0.67 on
planted ground-truth violations. The baseline exhaustive scanner
achieved recall of 1.00 but precision of only 0.71 --- it flagged
false-positive traps that the context-aware system correctly ignored.
Speedups ranged from 3x to 165x over exhaustive scanning. The
architecture is model-agnostic: the MCP server's semantic inference
layer can target local LLMs (Ollama, vLLM) or cloud endpoints, making
deployment viable across regulatory regimes with data residency
requirements.

\begin{center}\rule{0.5\linewidth}{0.5pt}\end{center}

\subsection{1. Introduction}\label{introduction}

Modern enterprises deploy compliance scanning tools to detect sensitive
data across endpoints and servers. These tools examine every file for
patterns matching regulated data: credit card numbers, API keys, social
security numbers, database credentials, and personally identifiable
information (PII). The approach is effective in principle but fails at
scale and fails in context.

A developer workstation might contain 30,000 files, of which fewer than
500 realistically contain anything worth scanning. The remaining scans
contribute nothing but computational overhead and, critically, expose
benign personal data to automated inspection engines. This creates a
paradox at the heart of compliance tooling: to enforce data privacy, the
tool inspects everything --- violating the very data minimisation
principles it claims to protect.

The second failure is contextual. A file called
\texttt{test\_card\_examples.md} inside a \texttt{/docs} directory
carries a fundamentally different risk profile from a file called
\texttt{.env.production} inside \texttt{/config}. Pattern-matching tools
flag both identically. False-positive rates above 80\% are common in
production deployments {[}1{]}, forcing security teams to manually
triage thousands of alerts before finding anything actionable.

Both failures share a root cause: current tools have no understanding of
\emph{where} a file sits in the context of the overall system. They
cannot distinguish test fixtures from production configurations,
documentation examples from real credentials, or sample data from
customer exports.

The Model Context Protocol (MCP) {[}2{]} creates a new architectural
possibility. MCP defines a standard interface through which language
models interact with external tools via structured JSON-RPC calls. An
MCP server can expose filesystem operations --- listing directories,
reading metadata, computing risk scores --- as individual tools that an
LLM orchestrates through a client. This structured tool boundary enables
a clean separation: the system can analyze filesystem \emph{structure}
through metadata-only MCP tools without ever invoking content-reading
tools.

This paper makes four contributions:

\begin{enumerate}
\def\labelenumi{\arabic{enumi}.}
\item
  \textbf{A three-phase MCP-based architecture} that separates
  structural analysis (Phase 1), risk profiling (Phase 2), and targeted
  content scanning (Phase 3), with file contents accessed only in the
  final phase.
\item
  \textbf{A working implementation} as a FastAPI server converted to an
  MCP server via \texttt{fastapi\_mcp}, exposing 15 tools for filesystem
  inspection, graph-based risk scoring, candidate selection, targeted
  scanning, and analyst feedback with adaptive weight learning.
\item
  \textbf{Empirical evaluation across six real repositories} (235 to
  29,339 files, including both open-source projects and private
  codebases) demonstrating 95\% scan reduction, 1.00 precision, and 0.67
  mean recall on planted ground-truth violations.
\item
  \textbf{A model-agnostic deployment framework} where the MCP server's
  semantic inference layer can target local LLMs for data residency
  compliance or cloud endpoints for convenience, with no architectural
  changes.
\end{enumerate}

\begin{center}\rule{0.5\linewidth}{0.5pt}\end{center}

\subsection{2. Background and Related
Work}\label{background-and-related-work}

\subsubsection{2.1 Data Loss Prevention
Systems}\label{data-loss-prevention-systems}

Commercial DLP systems --- Symantec DLP, Microsoft Purview, Digital
Guardian, Forcepoint --- operate by exhaustively scanning file contents
across endpoints, email, and cloud storage. These systems rely on
pattern matching (regular expressions for credit card numbers, social
security numbers), exact data matching (fingerprinting known sensitive
documents), and machine learning classifiers trained on labeled datasets
{[}3{]}. All approaches require reading file contents, and all scan
every file indiscriminately.

The Ponemon Institute reports that organizations spend an average of 395
hours per week managing DLP alerts, with false-positive rates between
50--80\% depending on the sensitivity of the detection rules {[}4{]}.
Gartner's 2024 Market Guide for DLP notes that ``organizations continue
to struggle with the operational overhead of managing DLP policies that
generate excessive false positives'' {[}5{]}.

\subsubsection{2.2 Model Context Protocol}\label{model-context-protocol}

The Model Context Protocol (MCP), introduced by Anthropic in November
2024 {[}2{]}, defines a client-server architecture where language models
interact with external systems through structured tool calls. An MCP
server exposes capabilities as tools with JSON Schema-defined inputs and
outputs. An MCP client (typically an LLM application) discovers
available tools and invokes them based on user intent.

MCP is transport-agnostic, supporting both stdio-based local servers and
HTTP-based remote servers with Server-Sent Events (SSE) for streaming.
The protocol has been adopted by multiple LLM platforms including Claude
Code, Cursor, Windsurf, and the OpenAI Agents SDK {[}6{]}.

For filesystem access, MCP servers can expose directory listing, file
reading, and metadata retrieval as individual tools. This creates a
natural enforcement boundary: an LLM can be given metadata-only tools
(list directories, read file stats) without content-reading tools,
structurally preventing it from accessing file contents during early
analysis phases.

\subsubsection{2.3 LLM-Based Security
Analysis}\label{llm-based-security-analysis}

Recent work has applied language models to security tasks including
vulnerability detection {[}7{]}, code review {[}8{]}, and threat
modeling {[}9{]}. Most approaches operate on file contents --- feeding
source code to the model for analysis. Our work differs in that the LLM
operates primarily on filesystem \emph{metadata} (paths, names,
extensions, permissions, sizes) rather than file contents, using
structural signals to identify which files warrant content inspection.

\subsubsection{2.4 Privacy-by-Design in Compliance
Tooling}\label{privacy-by-design-in-compliance-tooling}

The concept of Privacy by Design, formalized by Cavoukian {[}10{]} and
codified in GDPR Article 25, requires that data protection be embedded
into system design rather than added as an afterthought. Current DLP
architectures violate this principle: they access all data to determine
which data is sensitive. Hoepman {[}11{]} identifies data minimisation
as a core privacy design strategy --- systems should process only the
minimum data necessary for their purpose. Our architecture
operationalizes this principle by using structural metadata to minimize
the volume of file contents accessed during compliance scanning.

\begin{center}\rule{0.5\linewidth}{0.5pt}\end{center}

\subsection{3. System Architecture}\label{system-architecture}

The system is implemented as a FastAPI web application that is converted
into an MCP server using the \texttt{fastapi\_mcp} library {[}12{]}.
This means every FastAPI route automatically becomes an MCP tool,
callable by any MCP-compatible client. The server exposes 15 tools
organized into four categories.

\subsubsection{3.1 MCP Tool Surface}\label{mcp-tool-surface}

\textbf{Filesystem inspection tools} (no content access): -
\texttt{list\_directory} --- returns directory entries with names,
paths, and types - \texttt{get\_file\_metadata} --- returns file size,
permissions, owner, timestamps via \texttt{stat} -
\texttt{get\_directory\_structure} --- returns recursive tree structure
via \texttt{tree\ -J} with configurable depth -
\texttt{get\_disk\_usage} --- returns human-readable size via
\texttt{du}

\textbf{Intelligence and risk scoring tools} (no content access): -
\texttt{infer\_directory\_purpose} --- LLM-based semantic inference from
filenames and structure - \texttt{compute\_risk\_score} --- returns
weighted risk score with factor breakdown - \texttt{explain\_risk} ---
returns natural language risk explanation -
\texttt{identify\_scan\_candidates} --- deterministic candidate ranking
by extension, keywords, size, permissions

\textbf{Scanning tools} (content access --- Phase 3 only): -
\texttt{scan\_file\_sensitive\_data} --- regex-based pattern detection
on a single file - \texttt{scan\_candidate\_files} --- batch scan of
selected candidates - \texttt{scan\_directory\_sensitive\_data} ---
exhaustive recursive scan (baseline mode) -
\texttt{validate\_compliance} --- classifies scan results as COMPLIANT,
REVIEW, or NON\_COMPLIANT

\textbf{Feedback tools}: - \texttt{submit\_feedback} --- records analyst
TP/FP labels with notes - \texttt{update\_risk\_model} --- adjusts
scoring weights based on accumulated feedback -
\texttt{run\_full\_analysis} --- executes the complete three-phase
pipeline

The critical architectural property is the \textbf{tool boundary}:
filesystem inspection and risk scoring tools never read file contents.
Content access occurs only through scanning tools, which are invoked
only against files that survived Phase 2 filtering. An MCP client can
enforce this boundary by restricting which tools are available in each
phase.

\subsubsection{3.2 Phase 1: Structural
Discovery}\label{phase-1-structural-discovery}

Phase 1 collects filesystem metadata without reading any file contents.
The implementation uses Linux-native commands (\texttt{tree\ -J} for
recursive directory structure, \texttt{stat} for per-file metadata) with
Python fallbacks for cross-platform compatibility.

The collected metadata is assembled into a \texttt{networkx} graph
where: - Nodes represent files and directories, annotated with name,
path, depth, file type, size, permissions, and owner - Parent-child
edges represent directory containment - Name-similarity edges connect
files with similar stems and extensions, detected via bucket-based
prefix matching

This graph enables structural reasoning that goes beyond individual file
attributes. A \texttt{.env} file at the root of a project with high
degree centrality carries different risk than an \texttt{.env} file
nested inside a test fixture directory.

\subsubsection{3.3 Phase 2: Risk
Profiling}\label{phase-2-risk-profiling}

Phase 2 assigns a risk score \(R(f) \in [0, 1]\) to every file node in
the graph. The score is computed as a sigmoid-normalized weighted sum of
four signals:

\[R(f) = \sigma\left(\left(\sum_{i} w_i \cdot s_i(f)\right) - \theta\right)\]

where \(\sigma\) is the logistic sigmoid, \(\theta\) is a centering
threshold, and the signal components are:

\begin{itemize}
\tightlist
\item
  \textbf{Keyword signal} \(s_1(f)\): proportion of sensitive keywords
  (secret, token, password, credential, env, config, backup, export,
  key, pem, and 18 others) found in the filename, capped at 1.0
\item
  \textbf{Semantic signal} \(s_2(f)\): confidence and risk signals from
  LLM-based directory purpose inference
\item
  \textbf{Graph signal} \(s_3(f)\): degree centrality from the
  filesystem graph, elevated for nodes with high connectivity
\item
  \textbf{Metadata signal} \(s_4(f)\): composite of file extension risk
  (\texttt{.env}, \texttt{.pem}, \texttt{.csv}, \texttt{.sql} score
  higher), permission exposure (world-readable files), and file size
  (very large or very small files penalized)
\end{itemize}

Weights \(w_i\) default to \([0.35, 0.25, 0.20, 0.20]\) for keyword,
semantic, graph, and metadata respectively. These weights are adaptive:
the feedback loop (Section 3.6) adjusts them based on accumulated
true-positive and false-positive labels from analyst review.

Files are ranked by \(R(f)\) and the top 1--5\% are selected for Phase 3
scanning. The selection bounds are configurable
(\texttt{min\_scan\_percent}, \texttt{max\_scan\_percent}).

\textbf{Candidate selector (alternative path).} In addition to
graph-based scoring, the system provides a deterministic candidate
selector that operates without LLM calls. It classifies files by
extension (high-value text, source code, binary/media, parser-needed),
applies keyword scoring against mode-specific keyword sets (credentials,
PII, or broad), and incorporates file size, permission, and hidden-file
signals. This path is faster and fully reproducible, trading semantic
understanding for determinism.

\subsubsection{3.4 Phase 3: Targeted Content
Scanning}\label{phase-3-targeted-content-scanning}

Only files selected in Phase 2 undergo content inspection. The scanner
reads files asynchronously in configurable chunks (default 8,192 bytes)
to avoid loading entire files into memory. Pattern matching uses
compiled regular expressions for:

\begin{itemize}
\tightlist
\item
  \textbf{Email addresses}: standard RFC-compliant pattern
\item
  \textbf{Payment card numbers (PAN)}: 13--19 digit sequences
\item
  \textbf{API keys}: \texttt{sk\_}, \texttt{rk\_}, \texttt{pk\_}
  prefixed tokens
\item
  \textbf{AWS access keys}: \texttt{AKIA} prefix followed by 16
  alphanumeric characters
\item
  \textbf{Indian PAN IDs}: \texttt{AAAAA9999A} format (5 letters, 4
  digits, 1 letter)
\item
  \textbf{Aadhaar numbers}: 12-digit numbers starting with 2--9
\item
  \textbf{Credential keywords}: password, secret, api\_key, private key,
  ssn, social security
\end{itemize}

Each scan result is passed through compliance validation, which
classifies findings by severity. Files with API key, AWS key, PAN,
Indian PAN ID, or Aadhaar matches are classified as NON\_COMPLIANT.
Files with keyword-only matches are classified as REVIEW. Files with no
detections are COMPLIANT.

\subsubsection{3.5 Agent Orchestration}\label{agent-orchestration}

The MCP server is designed to be consumed by any MCP-compatible client.
Our reference client uses \texttt{mcp-use} {[}13{]} with a
\texttt{langchain}-based LLM (configurable --- tested with OpenRouter
endpoints and local Ollama models). The client loads a system prompt
that instructs the LLM to:

\begin{enumerate}
\def\labelenumi{\arabic{enumi}.}
\tightlist
\item
  Prefer selective workflows over brute-force scans
\item
  Call \texttt{identify\_scan\_candidates} before scanning a directory
\item
  Use \texttt{scan\_candidate\_files} on ranked candidates rather than
  recursive scanning
\item
  Explain whether results came from selective or exhaustive scanning
\end{enumerate}

This architecture means the LLM acts as an orchestration layer that
decides \emph{which} MCP tools to invoke and in what sequence. The
server remains a stateless tool provider.

\subsubsection{3.6 Feedback Loop}\label{feedback-loop}

Analyst feedback is persisted in a SQLite database as path-label pairs
(TP or FP) with optional notes. The \texttt{update\_risk\_model} tool
recomputes scoring weights based on accumulated feedback:

\begin{itemize}
\tightlist
\item
  High true-positive ratio increases keyword and semantic weights (the
  system is correctly identifying risky files, so double down on those
  signals)
\item
  High false-positive ratio increases graph and metadata weights
  (keyword/semantic signals are noisy, rely more on structural signals)
\end{itemize}

Weights are normalized to sum to 1.0 after adjustment. This creates a
per-deployment learning loop where the system adapts to the specific
filesystem patterns of each environment.

\begin{center}\rule{0.5\linewidth}{0.5pt}\end{center}

\subsection{4. Evaluation}\label{evaluation}

\subsubsection{4.1 Methodology}\label{methodology}

We evaluated the system against six real repositories spanning different
sizes, languages, and project types:

{\def\LTcaptype{none} % do not increment counter
\begin{longtable}[]{@{}llrrl@{}}
\toprule\noalign{}
Repository & Source & Files & Directories & Primary Language \\
\midrule\noalign{}
\endhead
\bottomrule\noalign{}
\endlastfoot
Flask & Open-source (GitHub) & 235 & 51 & Python \\
Django & Open-source (GitHub) & 6,944 & 3,265 & Python \\
Kshetra & Private & 505 & 115 & Mixed \\
veriflow & Private & 523 & 173 & Mixed \\
Yukthi & Private & 4,150 & 643 & Mixed \\
fairmind & Private & 29,339 & 1,229 & TypeScript/Python \\
\end{longtable}
}

\textbf{Ground truth.} For each repository, we planted five violation
files and two false-positive traps inside a
\texttt{.eval\_ground\_truth/} subdirectory:

\emph{Violation files (should be detected):} 1. \texttt{.env.production}
--- database credentials, AWS keys, Stripe live key, JWT secret 2.
\texttt{customer\_export.csv} --- customer PII with emails, phone
numbers, Indian PAN IDs 3. \texttt{db\_migrate.sh} --- shell script with
embedded database password 4. \texttt{payment\_handler.py} --- Python
file with hardcoded credit card number in a debug comment 5.
\texttt{employees\_q1\_2026.json} --- employee records with Aadhaar
numbers and salary data

\emph{False-positive traps (should NOT be detected):} 1.
\texttt{test\_payment\_validation.py} --- industry-standard test card
numbers (4111111111111111, etc.) 2. \texttt{api\_guide.md} ---
documentation with \texttt{YOUR\_API\_KEY\_HERE} placeholder

\textbf{Systems compared.} Three scanning approaches were evaluated on
each repository:

\begin{itemize}
\tightlist
\item
  \textbf{System A (Baseline)}: Exhaustive brute-force scan of every
  file using the same regex scanner
\item
  \textbf{System B (Candidate Selector)}: Deterministic heuristic
  ranking followed by targeted scan of top 100 candidates
\item
  \textbf{System B2 (Graph Pipeline)}: Full three-phase pipeline with
  graph-based risk scoring, selecting top 5\% of files
\end{itemize}

\textbf{Metrics.} Scan reduction (percentage of files not scanned),
recall (proportion of ground-truth violations detected), precision
(proportion of detections that are true violations), F1 score, and
wall-clock execution time. All evaluations were run on a single macOS
machine (Darwin 25.2.0) with LLM semantic inference disabled for
reproducibility.

\subsubsection{4.2 Results}\label{results}

\textbf{Table 1: Graph Pipeline (System B2) --- Primary Results}

{\def\LTcaptype{none} % do not increment counter
\begin{longtable}[]{@{}
  >{\raggedright\arraybackslash}p{(\linewidth - 14\tabcolsep) * \real{0.1558}}
  >{\raggedleft\arraybackslash}p{(\linewidth - 14\tabcolsep) * \real{0.1429}}
  >{\raggedleft\arraybackslash}p{(\linewidth - 14\tabcolsep) * \real{0.1169}}
  >{\raggedleft\arraybackslash}p{(\linewidth - 14\tabcolsep) * \real{0.1429}}
  >{\raggedleft\arraybackslash}p{(\linewidth - 14\tabcolsep) * \real{0.1299}}
  >{\raggedleft\arraybackslash}p{(\linewidth - 14\tabcolsep) * \real{0.1039}}
  >{\raggedleft\arraybackslash}p{(\linewidth - 14\tabcolsep) * \real{0.1429}}
  >{\raggedleft\arraybackslash}p{(\linewidth - 14\tabcolsep) * \real{0.0649}}@{}}
\toprule\noalign{}
\begin{minipage}[b]{\linewidth}\raggedright
Repository
\end{minipage} & \begin{minipage}[b]{\linewidth}\raggedleft
Total Files
\end{minipage} & \begin{minipage}[b]{\linewidth}\raggedleft
Scanned
\end{minipage} & \begin{minipage}[b]{\linewidth}\raggedleft
Reduction
\end{minipage} & \begin{minipage}[b]{\linewidth}\raggedleft
Time (s)
\end{minipage} & \begin{minipage}[b]{\linewidth}\raggedleft
Recall
\end{minipage} & \begin{minipage}[b]{\linewidth}\raggedleft
Precision
\end{minipage} & \begin{minipage}[b]{\linewidth}\raggedleft
F1
\end{minipage} \\
\midrule\noalign{}
\endhead
\bottomrule\noalign{}
\endlastfoot
Flask & 235 & 12 & 94.9\% & 0.06 & 0.40 & 1.00 & 0.57 \\
Kshetra & 505 & 26 & 94.8\% & 0.24 & 0.80 & 1.00 & 0.89 \\
veriflow & 523 & 27 & 94.8\% & 0.24 & 0.40 & 1.00 & 0.57 \\
Yukthi & 4,150 & 208 & 95.0\% & 4.90 & 0.80 & 1.00 & 0.89 \\
Django & 6,944 & 348 & 95.0\% & 1.95 & 0.80 & 1.00 & 0.89 \\
fairmind & 29,339 & 1,467 & 95.0\% & 19.27 & 0.80 & 1.00 & 0.89 \\
\textbf{Mean} & & & \textbf{94.9\%} & & \textbf{0.67} & \textbf{1.00} &
\textbf{0.78} \\
\end{longtable}
}

\textbf{Table 2: Three-System Comparison}

{\def\LTcaptype{none} % do not increment counter
\begin{longtable}[]{@{}
  >{\raggedright\arraybackslash}p{(\linewidth - 6\tabcolsep) * \real{0.1194}}
  >{\raggedleft\arraybackslash}p{(\linewidth - 6\tabcolsep) * \real{0.2985}}
  >{\raggedleft\arraybackslash}p{(\linewidth - 6\tabcolsep) * \real{0.3134}}
  >{\raggedleft\arraybackslash}p{(\linewidth - 6\tabcolsep) * \real{0.2687}}@{}}
\toprule\noalign{}
\begin{minipage}[b]{\linewidth}\raggedright
Metric
\end{minipage} & \begin{minipage}[b]{\linewidth}\raggedleft
System A (Baseline)
\end{minipage} & \begin{minipage}[b]{\linewidth}\raggedleft
System B (Candidate)
\end{minipage} & \begin{minipage}[b]{\linewidth}\raggedleft
System B2 (Graph)
\end{minipage} \\
\midrule\noalign{}
\endhead
\bottomrule\noalign{}
\endlastfoot
Mean scan reduction & 0\% & 88.0\% & 94.9\% \\
Mean recall & 1.00 & 0.60 & 0.67 \\
Mean precision & 0.71 & 0.91 & 1.00 \\
Mean F1 & 0.83 & 0.71 & 0.78 \\
\end{longtable}
}

\textbf{Table 3: Speedup Over Baseline (System B2 vs System A)}

{\def\LTcaptype{none} % do not increment counter
\begin{longtable}[]{@{}lrrr@{}}
\toprule\noalign{}
Repository & Baseline Time (s) & Graph Time (s) & Speedup \\
\midrule\noalign{}
\endhead
\bottomrule\noalign{}
\endlastfoot
Flask & 0.26 & 0.06 & 4.3x \\
Kshetra & 38.78 & 0.24 & 162x \\
veriflow & 39.63 & 0.24 & 165x \\
Yukthi & 16.97 & 4.90 & 3.5x \\
Django & 5.05 & 1.95 & 2.6x \\
fairmind & 427.39 & 19.27 & 22x \\
\end{longtable}
}

\subsubsection{4.3 Analysis}\label{analysis}

\textbf{Scan reduction is consistent.} The graph pipeline maintains
94.8--95.0\% scan reduction across all six repositories regardless of
size (235 to 29,339 files). This consistency arises from the fixed
\texttt{max\_scan\_percent} parameter (5\%) in the pipeline
configuration.

\textbf{Precision is perfect on the selective scanner.} System B2
achieved 1.00 precision across all six repositories --- every file it
flagged as containing violations was a true violation. This is because
the graph pipeline's risk scoring naturally deprioritizes test files,
documentation, and other low-risk contexts. The baseline scanner (System
A) achieved only 0.71 precision because it flagged the false-positive
traps (test card numbers and placeholder credentials) as violations.

\textbf{Recall gap explained.} The graph pipeline's mean recall of 0.67
means it missed some planted violations. Analysis of the missed files
reveals a consistent pattern: \texttt{db\_migrate.sh} (a shell script)
and \texttt{payment\_handler.py} (a Python file) were missed in most
repositories because their filenames do not contain high-signal
keywords. The risk engine's keyword-based scoring assigns lower scores
to files with generic names like ``payment\_handler'' compared to
explicitly sensitive names like ``.env.production'' or
``customer\_export.csv''. This represents a fundamental tradeoff: the
system optimizes for precision at the cost of recall on files whose
names do not signal their sensitivity.

\textbf{Speedup scales with noise.} The largest speedups (162x, 165x)
occurred on repositories where the baseline scanner processed tens of
thousands of files (including \texttt{node\_modules} and other
dependency trees) while the graph pipeline's tree walker excluded these
directories. The speedup is modest (2.6--4.3x) on repositories where the
baseline file count is closer to the enumerated file count.

\begin{center}\rule{0.5\linewidth}{0.5pt}\end{center}

\subsection{5. Regulatory Mapping}\label{regulatory-mapping}

The architecture's design choices map directly to data protection
principles across jurisdictions.

\subsubsection{5.1 GDPR}\label{gdpr}

\textbf{Article 5(1)(c) --- Data Minimisation.} The 95\% scan reduction
directly operationalizes the requirement that personal data be
``adequate, relevant and limited to what is necessary.'' By analyzing
filesystem structure rather than file contents, the system accesses the
minimum data necessary to identify compliance-relevant files.

\textbf{Article 25 --- Data Protection by Design.} The MCP tool boundary
enforces privacy by design: metadata-only tools are structurally
separated from content-reading tools. The architecture does not merely
choose not to read files --- the Phase 1 and Phase 2 tools \emph{cannot}
read file contents.

\textbf{Article 32 --- Security of Processing.} Proactive identification
of high-risk files (credentials in production configurations,
unencrypted PII in export files) supports the requirement for
``appropriate technical and organisational measures'' to ensure
security.

\subsubsection{5.2 India's Digital Personal Data Protection Act
2023}\label{indias-digital-personal-data-protection-act-2023}

\textbf{Section 4 --- Purpose Limitation.} The system restricts content
inspection to contextually justified targets. Files are scanned only
when structural signals indicate compliance relevance, not as a blanket
surveillance measure.

\textbf{Section 8 --- Data Fiduciary Obligations.} The architecture
supports Data Fiduciaries' obligation to implement ``reasonable security
safeguards'' that are proportionate to the risk. Scanning 5\% of files
with 100\% precision is more proportionate than scanning 100\% of files
with 71\% precision.

\subsubsection{5.3 EU AI Act}\label{eu-ai-act}

The system's use of LLM-based semantic inference falls under the EU AI
Act's requirements for transparency and accuracy. The deterministic
candidate selector provides a fully explainable alternative path ---
every file's risk score can be decomposed into specific factors (keyword
matches, extension classification, permission signals) without relying
on opaque model inference. Organizations can choose the deterministic
path for auditability or the LLM-enhanced path for richer contextual
understanding.

\begin{center}\rule{0.5\linewidth}{0.5pt}\end{center}

\subsection{6. Discussion and
Limitations}\label{discussion-and-limitations}

\subsubsection{6.1 The Recall-Precision
Tradeoff}\label{the-recall-precision-tradeoff}

The system's mean recall of 0.67 versus the baseline's 1.00 is the
central tradeoff. We argue this is acceptable in practice for three
reasons. First, the missed files had generic names that did not signal
sensitivity --- in production, compliance-relevant files tend to have
more descriptive names (credentials, secrets, customer data exports).
Second, the system achieves perfect precision, meaning security teams
spend zero time triaging false positives. Third, the
\texttt{max\_scan\_percent} parameter is tunable: increasing it from 5\%
to 10\% would improve recall at the cost of scanning more files.

\subsubsection{6.2 Ground Truth
Limitations}\label{ground-truth-limitations}

Our ground truth is synthetic --- violations were deliberately planted
in known locations. This is standard practice in DLP evaluation {[}3{]}
but does not capture organically occurring violations whose placement
may differ from our assumptions. Future work should evaluate against
datasets of real compliance violations (with appropriate anonymization).

\subsubsection{6.3 Reproducibility}\label{reproducibility}

All evaluation runs were conducted with LLM semantic inference disabled
(\texttt{DLP\_MCP\_LLM\_ENABLED=false}). This means the results reflect
only the deterministic components: keyword matching, extension
classification, graph centrality, and metadata signals. Enabling the
semantic inference layer would likely improve recall by adding
contextual understanding of directory purposes, but at the cost of
non-deterministic behavior across runs.

\subsubsection{6.4 Platform Coverage}\label{platform-coverage}

Evaluation was conducted on a single macOS machine. The system includes
cross-platform compatibility (Python fallbacks for GNU \texttt{stat} and
\texttt{tree}), but performance characteristics may differ on Linux and
Windows endpoints. Cross-platform evaluation is planned as future work,
with the evaluation harness already supporting Linux and WSL.

\subsubsection{6.5 Adversarial Evasion}\label{adversarial-evasion}

The system has not been evaluated against adversarial scenarios where
sensitive data is deliberately placed in benign-looking locations (e.g.,
credentials stored in \texttt{readme.txt} or API keys embedded in image
EXIF metadata). An adversary with knowledge of the scoring heuristics
could craft filenames and directory structures to evade detection.
Defense against adversarial evasion is an open problem for any
metadata-based pre-filtering approach.

\begin{center}\rule{0.5\linewidth}{0.5pt}\end{center}

\subsection{7. Conclusion}\label{conclusion}

We presented a three-phase compliance scanning architecture that uses
the Model Context Protocol to separate structural analysis from content
inspection. The system achieves 95\% scan reduction with perfect
precision across six real repositories, demonstrating that filesystem
structure alone provides sufficient signal to identify the small
fraction of files that warrant compliance inspection.

The MCP-based architecture offers a structural advantage over monolithic
DLP systems: the tool boundary between metadata-only and content-reading
operations is enforced by the protocol itself, not by policy. This makes
data minimisation an architectural property rather than a configuration
choice.

The system is open-source and available at
https://github.com/LakshmanS27/novel\_mcp.

\begin{center}\rule{0.5\linewidth}{0.5pt}\end{center}

\subsection{References}\label{references}

{[}1{]} Ponemon Institute, ``The True Cost of Compliance with Data
Protection Regulations,'' 2024.

{[}2{]} Anthropic, ``Model Context Protocol Specification,''
https://modelcontextprotocol.io, 2024.

{[}3{]} M. Bishop, ``Computer Security: Art and Science,''
Addison-Wesley, 2018.

{[}4{]} Ponemon Institute, ``Data Loss Prevention: Managing Insider
Risk,'' 2023.

{[}5{]} Gartner, ``Market Guide for Data Loss Prevention,'' 2024.

{[}6{]} OpenAI, ``Agents SDK: MCP Support,''
https://openai.github.io/openai-agents-python/mcp/, 2025.

{[}7{]} B. Steenhoek et al., ``A Comprehensive Study of the Capabilities
of Large Language Models for Vulnerability Detection,''
arXiv:2403.17218, 2024.

{[}8{]} Z. Li et al., ``Automated Code Review with LLMs: A Survey,''
arXiv:2402.08301, 2024.

{[}9{]} A. Abdeen et al., ``LLM-Based Threat Modeling,'' IEEE S\&P
Workshops, 2024.

{[}10{]} A. Cavoukian, ``Privacy by Design: The 7 Foundational
Principles,'' Information and Privacy Commissioner of Ontario, 2009.

{[}11{]} J.-H. Hoepman, ``Privacy Design Strategies,'' IFIP SEC, 2014.

{[}12{]} fastapi-mcp, ``Convert FastAPI routes to MCP tools,''
https://github.com/tadata-org/fastapi-mcp, 2025.

{[}13{]} mcp-use, ``MCP client for LangChain agents,''
https://github.com/pietrozullo/mcp-use, 2025.

\end{document}
